


\section{Discussion}
\vheading{Taxonomy of CPU computation} \sysname models host-side overheads as wall-clock time measurements between API calls to the emulator. This allows arbitrary host logic to be abstracted away, while still accounting for the impact of these overheads on end-to-end latency. However, there are workloads where significant CPU computation is involved, and this could affect prediction accuracy if there are hardware differences between the machine used during emulation vs. the target cluster. This can be addressed by applying the per-operation prediction approach to CPU work instead of simply collecting a wall-clock time, though this may not be exposed through a narrow API surface like accelerators. A combination of these two models could enable more general CPU overhead estimates.

\vheading{Dynamic control flow} As a result of relying on emulation, \sysname does not model computation graphs where the control flow depends on the result of tensor computation. This assumption is shared with several DL compilers and parallelism search engines \cite{zheng2022alpa}, \cite{flexflow}. Mixture-of-Experts (MoE) architectures display this pattern --- while most expert-parallel kernels used for MoE training \cite{deepep2025, pplxrepo} remove the need for data-dependent control-flow, there are some implementations that use gating on the host.

For expert-parallel kernels, runtime predictors can be trained by encoding the input distribution during the profiling phase \cite{lin2025apexextensibledynamismawaresimulator, vidur}, keeping the rest of the \sysname flow unchanged. To handle host-side gating, annotations on the source model can be used to identify the gating function --- instead of returning a random tensor during emulation, we would sample a distribution to generate a spread of runtimes.

\vheading{SM Contention} \sysname assumes decoupling between network collectives and concurrent compute streams. As a result, while we are able to model overlapping streams and arbitrary synchronization, we cannot trivially model SM-level interference where network and compute kernels contend for resources. It could be possible to modify and extend the simulator to identify such patterns and scale predicted durations accordingly --- we leave this to future work.
